%\ifodd\value{page}\hbox{}\newpage\fi
\chapter*{Śrī Lalitā Sahasranāma — Śloka 11}
\addcontentsline{toc}{chapter}{Śloka 11 - Nomi 24,25}

\begin{sloka}
{\sanskritfont
\begin{tabular}{c}
निजसल्लापमाधुर्यविनिर्भर्त्सितकच्छपी ।\\
मन्दस्मितप्रभापूरमज्जत्कामेशमानसा ॥ ११॥
\end{tabular}

\vspace{1.5em}

{\middlesize \iastfont
nija-sallāpa-mādhurya-vinirbhartsita-kacchapī |\\
manda-smita-prabhā-pūra-majjat-kāmeśa-mānasā ||11||}

\vspace{1em}

{\middlesize
\anvaya{%
निजसल्लापमाधुर्यविनिर्भर्त्सितकच्छपी ।
मन्दस्मितप्रभापूरमज्जत्कामेशमानसा ।
}
}

\middlesize \iastfont \itmean{%
«Il cui parlare proprio, per la sua dolcezza, mette in ombra il suono della \textit{kacchapī};  
il cui sorriso lieve, colmo di splendore, immerge la mente di \textit{Kāmeśa}.»%
}

}
\end{sloka}

\wordmeaning{%
\wordline{निज-सल्लाप-माधुर्य-विनिर्भर्त्सित-कच्छपी}
{nija-sallāpa-mādhurya-\\vinirbhartsita-kacchapī}
{ il cui parlare proprio (\textit{nija-sallāpa}, parola originaria e veritiera), per la sua dolcezza (\textit{mādhurya}), mette in ombra (\textit{vinirbhartsita}) il suono della \textit{kacchapī}, la \textit{vīṇā} di \textit{Sarasvatī} (il cui corpo richiama la forma di una tartaruga), emblema del suono raffinato e tecnicamente perfetto; }%

\wordline{मन्द-स्मित-प्रभा-पूर-मज्जत्-कामेश-मानसा}
{manda-smita-prabhā-pūra-\\majjat-kāmeśa-mānasā}
{ il cui sorriso lieve (\textit{manda-smita}), come una piena di splendore (\textit{prabhā-pūra}), immerge (\textit{majjat}) la mente (\textit{mānasa}) di \textit{Kāmeśa}; }%
}


\textit{Śloka} 11 sposta l’attenzione dalla forma statica della bocca alla sua potenza espressiva. Il parlare di \textit{Lalitā} non è descritto come semplice suono, ma come parola propria e originaria, capace di superare persino la dolcezza della musica strumentale. La \textit{kacchapī}, la \textit{vīṇā} di \textit{Sarasvatī}, emblema del suono raffinato e costruito, risulta inadeguata di fronte a una parola che nasce dalla presenza cosciente.

La tradizione ricorda, in forme allusive, che persino \textit{Sarasvatī} cesserebbe di suonare, vinta dalla dolcezza della voce di \textit{Lalitā}, a indicare che la musica più alta è la parola consapevole della Dea. È una memoria che attraversa una tradizione contemplativa in cui il \textit{manda-smita} non è semplice abbellimento, ma forma di rivelazione.

Questa dolcezza non è un semplice effetto estetico. Il termine \textit{sallāpa} indica un parlare intimo, misurato, relazionale: una voce che non esibisce tecnica, ma comunica prossimità. Qui la parola è veicolo diretto della \textit{śakti}, non sua imitazione sonora.

Nel secondo \textit{pāda}, il sorriso introduce un livello ancora più sottile. Non è gesto seduttivo, ma campo luminoso: una “piena di splendore” in cui persino la mente di \textit{Kāmeśa} viene immersa. L’immagine non suggerisce dominio, bensì assorbimento spontaneo, attrazione senza costrizione.

Questo \textit{śloka} contiene due \textit{nāma}.  
\textbf{\textit{Nija-sallāpa-mādhurya-\\vinirbhartsita-kacchapī}} presenta la parola della Dea come dolcezza originaria che supera ogni suono artefatto.  
\textbf{\textit{Manda-smita-prabhā-pūra-majjat-kāmeśa-mānasā}} raffigura il sorriso come irradiazione nella quale anche la coscienza di \textit{Kāmeśa} si dissolve senza resistenza.
